\section{Equa��o de Calor e Grupos de Transforma��es}


\subsection{Grupos de Transforma��o}

Para a equa��o de calor mostraremos que esta equa��o permanece invariante sobre as seguintes transforma��es:
%
%
\begin{eqnarray}
u_{t} - u_{xx} &=& 0 \qquad \mbox{(Equa��o de Calor)} \label{calor}\\
t^* &=& t + \e \qquad \mbox{(Transla��es no tempo)} \label{tempo}\\
x^* &=& x + \e \qquad \mbox{(Transla��es em coordenadas espaciais)} \label{espacial}\\
u^* &=& u + \e \varphi(t,x) \qquad \mbox{(Princ�pio da Superposi��o)} \label{superposicao}\\
u^* &=& ue^{\e} 
\end{eqnarray}
%
%
Na equa��o \ref{superposicao} temos que $\varphi(t,x)$ � uma solu��o qualquer da equa��o \ref{calor}. Utilizando as rela��es \ref{tempo} e \ref{espacial}, e a regra da cadeia para fun��es de v�rias var�veis, tem-se que:
%
%
\begin{eqnarray}
\deli{}{t} &=& \deli{t^*}{t}\deli{}{t^*} +  \deli{x^*}{t}\deli{}{x^*} = \deli{}{t^*} \label{grupotempo}\\
\deli{}{x} &=& \deli{t^*}{x}\deli{}{t^*} +  \deli{x^*}{x}\deli{}{x^*} = \deli{}{x^*}\label{grupoespaco}
\end{eqnarray}
%
%
De forma que a equa��o de calor permane�a invariante sobre transla��es no tempo e nas coordenadas espaciais. Ou seja, tem-se que:
%
%
\[u_{t} - u_{xx} = u_{t^*} - u_{x^*x^*} = 0\]
%
%
Utilizando as rela��es \ref{grupotempo} e \ref{grupoespaco}, e a transforma��o $u^* = u + \e \varphi(t,x)$, encontra-se que:
%
%
\begin{eqnarray}
\nonumber (\partial_{t^*} - \partial_{x^*x^*})(u + \e \varphi(t,x)) &=& (\partial_{t} - \partial_{xx})(u + \e \varphi(t,x))\\
\nonumber &=& (\partial_{t} - \partial_{xx})u + (\partial_{t} - \partial_{xx})(\e \varphi(t,x))\\
\nonumber &=& (u_t - u_{xx} + \e(\varphi_t - \varphi_{xx})\\
\nonumber &=& 0
\end{eqnarray}
%
%
Da mesma forma, procedemos para a transforma��o $u^* = u e^\e$:
%
%
\begin{eqnarray}
\nonumber (\partial_{t^*} - \partial_{x^*x^*})(u e^\e) &=& (\partial_{t} - \partial_{xx})(u e^\e)\\
\nonumber &=& e^\e(u_t - u_{xx})\\
\nonumber &=& 0
\end{eqnarray}